Exception processing covers synchronous restartable CPU exceptions, traps, NMIs, external interrupts, and the
supervisor-entry state needed to return from them. The processor records a restartable context for CPU exceptions;
the operating system decides whether that context is retried, emulated, reported, or terminated. ICR selects the
interrupt vector table page and controls interrupt nesting; each IVT entry selects a handler and one of the
supervisor stack banks.

\subsection{Interrupt Vector Table}
The interrupt vector table occupies one 4096-byte page selected by \texttt{ICR.ivt\_page}. It contains 256 fixed-size
entries, each 16 bytes. The vector number selects the entry at \texttt{ICR.ivt\_page * 4096 + vector * 16}.

\begin{manuallistedformatdiagram}{Interrupt Vector Table Entry}{3}
\manualformatrowrange{entry bytes}{15}{0}{%
\manualformatfield{reserved}{7}
\manualformatfield{control}{1}
\manualformatfield{handler address}{8}
}
\end{manuallistedformatdiagram}


Bytes 0 through 7 contain the 64-bit interrupt handler address. Byte 8 is the entry control byte.
HP is the handler-present bit; if it is clear, delivery immediately becomes \texttt{DOUBLE\_FAULT}. SN selects one of
four interrupt stack banks, SSS0:SSP0 through SSS3:SSP3. Byte 8 bit 1 and bits 7 through 4 are reserved,
and bytes 9 through 15 are reserved.

\begin{manuallistedformatdiagram}{IVT Entry Control Byte}{2}
\manualformatrowrange{byte 8}{7}{0}{%
\manualformatfield{reserved}{4}
\manualformatfield{SN}{2}
\manualformatfield{0}{1}
\manualformatfield{HP}{1}
}
\end{manuallistedformatdiagram}


@IVT_ENTRY_TABLE@

\Needspace{3.4in}
\subsection{Vector Assignment}
@VECTOR_ASSIGNMENT_POLICY@

\texttt{SYSCALL} does not allocate or consume an IVT vector. It enters supervisor mode through the @SYSCALL_ENTRY_PATH@.

@VECTOR_RANGE_TABLE@

@CPU_VECTOR_TABLE@

\subsection{Entry and Return Rules}
The processor changes only the privilege bit on entry: STATUS.PM is set to supervisor mode. It also saves the
interrupted DBANK selector and sets DBANK to zero before fetching the handler. Interrupts are not automatically
masked merely because an exception or interrupt was taken. The interrupt stack frame is saved atomically. When
returning, the saved STATUS image is restored exactly as recorded in the frame, and the saved DBANK selector is
restored from the frame.

@EXCEPTION_ENTRY_TABLE@

\clearpage
\subsection{Supervisor Entry Stack Frame}
The fixed supervisor-entry frame is @BASE_FRAME_SIZE@ bytes: @FIXED_SLOT_COUNT@ 64-bit slots.
\texttt{FRAME\_CONTROL.frame\_size\_units} records the total frame size in @FRAME_SIZE_UNIT_BYTES@-byte units, so the base frame
records @BASE_FRAME_UNITS@ units before any optional payload. Offsets are measured from the saved stack pointer.

@STACK_FRAME_FIGURE@

@STACK_FRAME_TABLE@

\texttt{FRAME\_CONTROL.frame\_type} documents the stack frame format, M68000-style. Payload, if any, is appended after the
fixed frame in whole @PAYLOAD_BLOCK_SIZE@-byte blocks. \texttt{FRAME\_CONTROL.frame\_size\_units} gives the total number of
@FRAME_SIZE_UNIT_BYTES@-byte slots, and the frame type code determines how software interprets the allocated payload blocks.
IRET does not validate the frame type.

@FRAME_TYPE_TABLE@

\texttt{FRAME\_CONTROL} packs the saved DBANK selector, FLAGS image, and STATUS image with the frame metadata.
The fixed frame also reserves two 64-bit extension slots for future architected base-frame state.

@PAYLOAD_SLOT_TABLE@

\begin{manuallistedformatdiagram}{FRAME\_CONTROL Format}{3}
\manualformatrowrange{FRAME\_CONTROL[63:32]}{63}{32}{%
\manualformatfield{STATUS}{16}
\manualformatfield{FLAGS}{16}
}
\manualformatrowrange{FRAME\_CONTROL[31:0]}{31}{0}{%
\manualformatfield{reserved}{1}
\manualformatfield{DBANK}{4}
\manualformatfield{entry flags}{3}
\manualformatfield{frame type}{4}
\manualformatfield{idepth}{4}
\manualformatfield{frame size}{8}
\manualformatfield{vector}{8}
}
\end{manuallistedformatdiagram}
\vspace{4pt}


@FRAME_CONTROL_TABLE@

\paragraph{Repeat Fault Continuation.}
When \texttt{FRAME\_CONTROL.rep\_fault} is set, the \texttt{FAULT\_AUX} payload records the active repeat context.
\texttt{FAULT\_AUX.repeat\_kind} distinguishes \texttt{REP}/\texttt{REPcc}, \texttt{REPG-general}, and \texttt{REPG-fast} continuations.
For \texttt{REPG}, the saved \texttt{PC} is the faulting grouped instruction, and the group start is reconstructed
as \texttt{group\_start\_pc = fault\_pc - body\_start\_delta\_bytes}. The auxiliary word also records the
\texttt{REPG} body byte count, so returning with IRET preserves the continuation state: execution resumes at the
saved faulting instruction under the active repeat context and eventually repeats or exits through the byte-counted
body. Software that wants to abandon or emulate the repeated operation may edit the frame before returning.

@REPEAT_FAULT_AUX_TABLE@

\subsection{Interrupt Reset State}
@RESET_STATE_TABLE@
